\section{Linear Elliptic Equations}\label{chap:linear-elliptic-equations}

\subsection{Definitions}\label{sec:elliptic-equation-definitions}

We will begin our study of linear elliptic PDEs by considering differential operators of the form $L(\bm{x},D)u=0$ defined in \eqref{eq:general-linear-differential-expression}. Note that by Definition~\ref{def:higher-order-pde-type} a differential operator of order $m$ is elliptic at $\bm{x}_0$ if and only if
\begin{equation}\label{eq:elliptic-opening-principal-symbol-nonzero}
L^p(\bm{x}_0,\bm{\xi})=\sum_{|\alpha|=m}a_\alpha(\bm{x})\bm{\xi}^\alpha\ne0\quad\text{for every }\bm{\xi}\in\mathbb{R}^n\backslash\{0\}.
\end{equation}
In fact, we can show the following.

\begin{Lemma}
\label{lem:elliptic-operator-even-order-fixed-sign}
If a linear partial differential operator $L$ of order $m$ is elliptic at $\bm{x}_0\in\mathbb{R}^n$, $n>1$, then $m$ is an even integer ($m=2k$) and $\bm{\xi}\mapsto L^p(\bm{x}_0,\bm{\xi})$ takes on only one sign on $\bm{\xi}\ne0$.
\end{Lemma}

\begin{Proof}
By definition, $\bm{\xi}\mapsto L^p(\bm{x}_0,\bm{\xi})$ is continuous and takes on the value $0$ only at $\bm{\xi}=0$. Suppose $L^p(\bm{x}_0,\bm{\xi}_1)<0$ and $L^p(\bm{x}_0,\bm{\xi}_2)>0$, and then connect $\bm{\xi}_1$ and $\bm{\xi}_2$ using a path not going through $0$. As noted, $L^p(\bm{x}_0,\bm{\xi})$ must vary continuously along the path, taking on the value $0$. This is a contradiction.

It now follows that, for any $\bm{\xi}\in\mathbb{R}^n$,\[L^p(\bm{x}_0,\bm{\xi})\quad\text{and}\quad L^p(\bm{x}_0,-\bm{\xi})=(-1)^mL^p(\bm{x}_0,\bm{\xi})\]must have the same sign. This implies that $m$ is even.
\end{Proof}

In light of this result, we will use the following somewhat restricted definition of an elliptic operator for the remainder of the chapter.

\begin{Definition}
\label{def:positive-sign-elliptic-and-uniform-elliptic}
Let $\Omega\subseteq\mathbb{R}^n$ be a domain. We say that a linear partial differential operator
\begin{equation}\label{eq:even-order-elliptic-differential-operator}
L(\bm{x},D)=\sum_{|\alpha|\leq2k}a_\alpha(\bm{x})D^\alpha
\end{equation}
is \textbf{elliptic} in $\Omega$ if
\begin{equation}\label{eq:positive-sign-ellipticity-condition}
(-1)^k\sum_{|\alpha|=2k}a_\alpha(\bm{x})\bm{\xi}^\alpha>0\quad\text{for every }\bm{x}\in\Omega,\ \bm{\xi}\in\mathbb{R}^n\backslash\{0\}.
\end{equation}
We say that $L$ is \textbf{uniformly elliptic} in $\Omega$ if there exists a constant $\theta>0$ such that
\begin{equation}\label{eq:uniform-ellipticity-principal-symbol-bound}
(-1)^k\sum_{|\alpha|=2k}a_\alpha(\bm{x})\bm{\xi}^\alpha\geq\theta|\bm{\xi}|^{2k}\quad\text{for every }\bm{x}\in\Omega,\ \bm{\xi}\in\mathbb{R}^n\backslash\{0\}.
\end{equation}

\end{Definition}

\begin{Example}
\label{ex:laplacian-biharmonic-uniform-ellipticity}
The reader should recall the calculations of Chapter~\ref{chap:characteristics} which showed that the negative of the Laplacian $-\Delta$ (which is of order $2$) and the Biharmonic $\Delta^2$ (order $4$) are uniformly elliptic with $\theta=1$.
\end{Example}

\begin{Example}
\label{ex:second-order-negative-matrix-uniform-ellipticity}
A second-order operator in $n$ space dimensions of the form
\begin{equation}\label{eq:second-order-elliptic-nondivergence-operator}
Lu=a_{ij}(\bm{x})\frac{\partial^2u}{\partial x_i\partial x_j}+b_i(\bm{x})\frac{\partial u}{\partial x_i}+c(\bm{x})u
\end{equation}
is uniformly elliptic on a domain $\Omega$ provided there exists a constant $\theta$ such that
\begin{equation}\label{eq:second-order-ellipticity-negative-coefficient-matrix}
\bm{\xi}^T\bm{A}(\bm{x})\bm{\xi}>\theta|\bm{\xi}|^2
\end{equation}
for every $\bm{x}\in\Omega$. Here $\bm{A}(\bm{x})$ is the $n\times n$ matrix with components $-a_{ij}(\bm{x})$.
\end{Example}

In our discussion of existence and regularity theory below, it is convenient to put our differential operators in a form which is amenable to integration by parts.

\begin{Definition}
\label{def:general-divergence-form-operator}
We say that an operator is in \textbf{divergence form} if there are functions $a_{\sigma\gamma}:\Omega\to\mathbb{R}$ such that
\begin{equation}\label{eq:general-even-order-divergence-form-definition}
L(\bm{x},D)u=\sum_{0\leq|\sigma|,|\gamma|\leq k}(-1)^{|\sigma|}D^\sigma(a_{\sigma\gamma}(\bm{x})D^\gamma u).
\end{equation}

\end{Definition}

\begin{Remark}
\label{rem:divergence-form-ellipticity-symbol}
Note that an operator in divergence form is elliptic if and only if
\begin{equation}\label{eq:divergence-form-ellipticity-symbol-condition}
\sum_{|\sigma|,|\gamma|=k}\bm{\xi}^\sigma a_{\sigma\gamma}(\bm{x})\bm{\xi}^\gamma>0\quad\text{for every }\bm{x}\in\Omega,\ \bm{\xi}\in\mathbb{R}^n\backslash\{0\},
\end{equation}
and uniformly elliptic if and only if there exists $\theta>0$ such that
\begin{equation}\label{eq:divergence-form-uniform-ellipticity-condition}
\sum_{|\sigma|,|\gamma|=k}\bm{\xi}^\sigma a_{\sigma\gamma}(\bm{x})\bm{\xi}^\gamma>\theta|\bm{\xi}|^{2k}\quad\text{for every }\bm{x}\in\Omega,\ \bm{\xi}\in\mathbb{R}^n\backslash\{0\}.
\end{equation}

\end{Remark}
